\chapter{专题——祷告}
\section{什么是祷告}


\subsection{長大痲瘋的祷告 太8:1-4 learned}
耶穌下了山，有許多人跟著他。 2 有一個長大痲瘋的來拜他，說：「主若肯，必能叫我潔淨了。」 3 耶穌伸手摸他，說：「我肯，你潔淨了吧！」他的大痲瘋立刻就潔淨了。耶穌對他說：「你切不可告訴人，只要去把身體給祭司察看，獻上摩西所吩咐的禮物，對眾人做證據。」

\subsection{耶稣教导门徒的祷告 太6:9-14 learnning}
所以，你們禱告要這樣說：『我們在天上的父，願人都尊你的名為聖， 10 願你的國降臨，願你的旨意行在地上，如同行在天上。 11 我們日用的飲食，今日賜給我們。 12 免我們的債，如同我們免了人的債。 13 不叫我們遇見試探，救我們脫離凶惡（脫離惡者）。因為國度、權柄、榮耀，全是你的，直到永遠。阿們。』 14 你們饒恕人的過犯，你們的天父也必饒恕你們的過犯。 15 你們不饒恕人的過犯，你們的天父也必不饒恕你們的過犯。


\subsection{彼得在海面沉水前的祷告 太14:22-33}
耶穌隨即催門徒上船，先渡到那邊去，等他叫眾人散開。 23 散了眾人以後，他就獨自上山去禱告。到了晚上，只有他一人在那裡。 24 那時船在海中，因風不順，被浪搖撼。 25 夜裡四更天，耶穌在海面上走，往門徒那裡去。 26 門徒看見他在海面上走，就驚慌了，說：「是個鬼怪！」便害怕，喊叫起來。 27 耶穌連忙對他們說：「你們放心！是我，不要怕！」 28 彼得說：「主，如果是你，請叫我從水面上走到你那裡去。」 29 耶穌說：「你來吧！」彼得就從船上下去，在水面上走，要到耶穌那裡去。 30 只因見風甚大，就害怕，將要沉下去，便喊著說：「主啊，救我！」 31 耶穌趕緊伸手拉住他，說：「你這小信的人哪，為什麼疑惑呢？」 32 他們上了船，風就住了。 33 在船上的人都拜他，說：「你真是神的兒子了！」

\subsection{門徒求耶稣平靜風和海的祷告 路8:22-25}
有一天，耶穌和門徒上了船，對門徒說：「我們可以渡到湖那邊去。」他們就開了船。 23 正行的時候，耶穌睡著了。湖上忽然起了暴風，船將滿了水，甚是危險。 24 門徒來叫醒了他，說：「夫子！夫子！我們喪命啦！」耶穌醒了，斥責那狂風大浪，風浪就止住，平靜了。 25 耶穌對他們說：「你們的信心在哪裡呢？」他們又懼怕又稀奇，彼此說：「這到底是誰？他吩咐風和水，連風和水也聽從他了！」


\subsection{百夫长为仆人代求的祷告 太8:5-13}
耶穌進了迦百農，有一個百夫長進前來，求他說： 6 「主啊，我的僕人害癱瘓病，躺在家裡甚是疼苦。」 7 耶穌說：「我去醫治他。」 8 百夫長回答說：「主啊，你到我舍下我不敢當，只要你說一句話，我的僕人就必好了。 9 因為我在人的權下，也有兵在我以下，對這個說『去！』他就去，對那個說『來！』他就來，對我的僕人說『你做這事！』他就去做。」 10 耶穌聽見就稀奇，對跟從的人說：「我實在告訴你們：這麼大的信心，就是在以色列中我也沒有遇見過！ 11 我又告訴你們：從東從西，將有許多人來，在天國裡與亞伯拉罕、以撒、雅各一同坐席； 12 唯有本國的子民，竟被趕到外邊黑暗裡去，在那裡必要哀哭切齒了。」 13 耶穌對百夫長說：「你回去吧！照你的信心給你成全了。」那時，他的僕人就好了。

\subsection{父亲求被鬼附的儿子得医治的祷告 learned}
\subsubsection{可9:14-27}
耶穌到了門徒那裡，看見有許多人圍著他們，又有文士和他們辯論。
眾人一見耶穌，都甚希奇，就跑上去問他的安。
耶穌問他們說：你們和他們辯論的是什麼？
眾人中間有一個人回答說：\textcolor{red}{夫子，我帶了我的兒子到你這裡來，他被啞巴鬼附著。
無論在哪裡，鬼捉弄他，把他摔倒，他就口中流沫，咬牙切齒，身體枯乾。我請過你的門徒把鬼趕出去，他們卻是不能}。
耶穌說：噯！不信的世代啊，我在你們這裡要到幾時呢？我忍耐你們要到幾時呢？把他帶到我這裡來吧。
他們就帶了他來。他一見耶穌，鬼便叫他重重地抽瘋，倒在地上，翻來覆去，口中流沫。
耶穌問他父親說：他得這病有多少日子呢？回答說：\textcolor{red}{從小的時候。
鬼屢次把他扔在火裡、水裡，要滅他。你若能做什麼，求你憐憫我們，幫助我們}。
耶穌對他說：你若能信，在信的人，凡事都能。
孩子的父親立時喊著說（有古卷：立時流淚地喊著說）：\textcolor{red}{我信！但我信不足，求主幫助}。
耶穌看見眾人都跑上來，就斥責那污鬼說：你這聾啞的鬼，我吩咐你從他裡頭出來，再不要進去！
那鬼喊叫，使孩子大大的抽了一陣瘋，就出來了。孩子好像死了一般。以致眾人多半說：他是死了。
但耶穌拉著他的手，扶他起來，他就站起來了。

\subsubsection{问题}
1，有几个人物\\
2，父亲为什么来\\
3，"眾人中間有一個人回答"的“一個人”是谁？和被鬼附的孩子是什么关系？\\
4，在耶稣来之前，父亲带着孩子先求了谁？\\
5，都谁向耶稣做祷告了\\
6，父亲祷告的内容包括什么？\\
\hspace{1cm}   \textcolor{blue}{a，叙述病情}
\hspace{1cm}   \textcolor{blue}{b，和两次祈求}\\
7，哪个祷告蒙耶稣成就了？\\
8，从父亲的身上我们学习了什么？\\
\hspace{1cm}    \textcolor{blue}{a，有难处到主面前求}
\hspace{1cm}    \textcolor{blue}{b，恳切的求（流泪）}
\hspace{1cm}    \textcolor{blue}{c，照着神的心意求(我信！但我信不足，求主幫助)}

\section{为什么要祷告}


\section{祷告的种类}
\section{祷告的内容}

“耶和华啊，早晨你必听我的声音。早晨我必向你陈明我的心意，并要警醒”（诗篇5：3）。“神啊，我曾求告你，因为你必应允我：求你向我侧耳，听我的言语”（诗篇17：6）。

大衛作詩（詩55:17）說，一天要三次禱告神；大衛又作詩（詩119:164）說，一天要七次讚美神。

\section{祷告后的结果}
\subsection{同意}
\subsection{不同意}
\subsection{神的意见和我们的不一样}
\subsection{等待}
